% =====================================================================
%  CHAPTER 12: 儿童少年卫生服务和学校卫生监督（对应大纲1第十二章）
% =====================================================================
\chapter{儿童少年卫生服务和学校卫生监督}
\setcounter{chapter}{12}
\chapterintro{
\textbf{掌握：}儿童少年卫生服务的七大内容；健康促进学校的概念与标准；学校卫生监督的内容与程序。\newline
\textbf{熟悉：}儿童少年卫生服务体系（儿童保健+学校卫生）；中国健康促进学校的发展阶段。
}

% ----- 12.1 儿童少年卫生服务 -----
\section{儿童少年卫生服务}

\subsection{概念}

\begin{defbox}
\textbf{儿童少年卫生服务（health service for children and adolescents）：}以儿童少年为主要服务对象，通过政府、社会、学校和家庭的共同努力，为儿童少年提供系统、综合、连续的卫生保健服务，旨在保护和促进儿童少年的身心健康，预防和控制疾病与伤害，促进生长发育潜能的充分发挥，为终身健康奠定基础。

儿童少年卫生服务是儿童少年卫生学的重要组成部分，也是实现"健康中国"战略的关键环节。其服务内容涵盖从胎儿期到青年期各年龄段的身心健康需求，贯穿生命早期发展的全过程。
\end{defbox}

\subsection{卫生服务体系}

\begin{keybox}
\textbf{儿童少年卫生服务体系由两大体系构成：}

\begin{enumerate}[leftmargin=*]
    \item \imp{儿童保健服务体系}——主要覆盖0--6岁儿童，重点加强对3岁以下婴幼儿的照护服务和健康管理。由各级妇幼保健机构、社区卫生服务中心、乡镇卫生院等组成，提供涵盖新生儿访视、儿童健康检查、生长发育监测、免疫规划、营养指导等系统保健服务。

    \item \imp{学校卫生服务体系}——主要覆盖6--24岁在校青少年儿童，由疾病预防控制中心（学校卫生科/所）、中小学卫生保健所、校医院/卫生室、社区卫生服务中心等组成，负责学生健康监测、常见病防控、健康教育、学校环境卫生监督等工作。
\end{enumerate}

\textbf{两大体系的衔接与协作：}儿童保健服务与学校卫生服务在学龄前期形成衔接，通过入学健康检查、预防接种证查验等机制实现信息共享和服务延续，共同构成覆盖0--24岁全年龄段的儿童少年卫生服务网络。
\end{keybox}

\begin{defbox}
\textbf{服务体系的核心特征：}
\begin{enumerate}[leftmargin=*]
    \item \imp{连续性}：从胎儿期到青年期，各年龄段服务无缝衔接
    \item \imp{综合性}：涵盖生长发育、疾病防控、心理健康、营养、安全等多维度
    \item \imp{多部门协作}：教育、卫生、体育、市场监管等多部门共同参与
    \item \imp{学校为基础}：学校是实施儿童少年卫生服务的主要场所
\end{enumerate}
\end{defbox}

\subsection{卫生服务内容及实施}

\begin{keybox}
\textbf{儿童少年卫生服务共有七大内容：}
\end{keybox}

\subsubsection{（一）生长发育与健康监测}

\begin{defbox}
\textbf{主要内容：}
\begin{enumerate}[leftmargin=*]
    \item \imp{生长发育监测}：定期对学生进行身高、体重、胸围等体格测量，评价生长发育水平、速度与匀称度。
    \item \imp{健康体检}：包括内外科检查、五官科检查、实验室检查等，建立学生健康档案。
    \item \imp{健康状况评估}：对学生个体和群体健康状况进行综合评价，提出相应干预建议。
    \item \imp{常见病监测}：重点监测视力不良、龋齿、营养不良、超重肥胖、脊柱弯曲异常等学生常见病。
    \item \imp{健康相关行为监测}：调查学生健康危险行为（吸烟、饮酒、不合理膳食、缺乏体力活动等）。
\end{enumerate}

\textbf{实施要求：}每学年进行一次常规健康体检，建立"一人一档"健康档案，实现学生健康状况的动态追踪管理。监测数据为国家制定学校卫生政策提供科学依据。
\end{defbox}

\subsubsection{（二）健康教育和健康促进}

\begin{keybox}
\textbf{实施途径：}
\begin{enumerate}[leftmargin=*]
    \item 健康教育课程（学科渗透+专题教育）
    \item 健康教育活动（主题班会、健康日/周/月活动）
    \item 健康咨询与指导（个体咨询、团体辅导）
    \item 校园健康文化建设（健康宣传栏、广播、新媒体等）
\end{enumerate}

\vspace{4pt}
\textbf{中小学健康教育内容5个领域：}
\begin{enumerate}[leftmargin=*]
    \item \imp{健康行为与生活方式}：培养学生良好的饮食、作息、运动、个人卫生等日常生活习惯。
    \item \imp{疾病预防}：传授常见传染病和慢性病的预防知识与技能。
    \item \imp{生长发育与青春期保健}：帮助学生正确认识自身生长发育过程，掌握青春期保健知识。
    \item \imp{心理健康}：培养积极的心理品质，提高情绪管理和人际交往能力。
    \item \imp{安全应急与避险}：学习常见意外伤害的预防与急救知识，提升安全意识和自救互救能力。
\end{enumerate}
\end{keybox}

\subsubsection{（三）常见病防控与慢性病管理}

\begin{keybox}
\textbf{"多病共防"理念：}学生常见病和慢性病往往共享相同的危险因素（不健康饮食、缺乏运动、熬夜等），应采取综合防控策略，实现多病共防、一举多得。

\textbf{常见病防控重点：}
\begin{enumerate}[leftmargin=*]
    \item \imp{视力不良防控}：改善教室采光照明、规范眼保健操、增加户外活动时间、"一尺一拳一寸"读写姿势
    \item \imp{龋齿防控}：窝沟封闭、氟化物应用、口腔健康教育、定期口腔检查
    \item \imp{营养不良防控}：营养教育、营养改善计划、生长发育监测与干预
    \item \imp{超重肥胖防控}：合理膳食指导、增加体力活动、减少静态行为
    \item \imp{脊柱弯曲异常防控}：正确坐姿教育、书包重量控制、脊柱健康筛查
\end{enumerate}

\textbf{慢性病管理：}对已患有哮喘、糖尿病、高血压等慢性病的学生，建立校医与专科医师联动管理机制，制定个体化健康管理计划，确保在校期间的健康安全。
\end{keybox}

\subsubsection{（四）学校心理卫生服务}

\begin{keybox}
\textbf{关键原则：}
\begin{enumerate}[leftmargin=*]
    \item 面向全体学生与关注特殊群体相结合
    \item 发展性辅导与预防性干预相结合
    \item 保密原则与生命安全相平衡
    \item 学校主导与家校社协同相结合
\end{enumerate}
\end{keybox}

\subsubsection{（五）学校营养服务}

\begin{keybox}
\textbf{学校营养服务实施内容：}
\begin{enumerate}[leftmargin=*]
    \item 营养健康教育课程和宣传
    \item 学生营养餐/学校供餐的营养指导
    \item 学生营养状况监测与评价
    \item 营养不良/营养过剩学生的个体化营养干预
    \item 食品安全管理（食堂卫生、饮用水安全等）
\end{enumerate}
\end{keybox}

\subsubsection{（六）校园安全与伤害暴力预防}

\begin{defbox}
\textbf{主要内容：}
\begin{enumerate}[leftmargin=*]
    \item \imp{校园设施安全}：校舍、运动场所、实验室等设施设备的定期安全检查和维护
    \item \imp{交通安全}：校车安全管理、上下学交通安全教育
    \item \imp{食品和饮用水安全}：学校食堂和饮水设施的卫生监管
    \item \imp{意外伤害预防}：溺水、跌落、烧烫伤、中毒等常见意外伤害的预防教育
    \item \imp{校园暴力与欺凌预防}：建立反校园欺凌制度和报告机制，开展同理心教育和冲突解决技能训练
    \item \imp{自然灾害与突发事件应急}：制定应急预案，定期开展应急疏散演练
    \item \imp{儿童保护}：预防虐待与忽视，建立侵害未成年人案件强制报告制度
\end{enumerate}
\end{defbox}

\subsubsection{（七）体育与体力活动促进}

\begin{keybox}
\textbf{实施保障：}
\begin{enumerate}[leftmargin=*]
    \item 体育场地和设施器材达标配置
    \item 体育教师配备与专业培训
    \item 学生体质健康标准测试与评价
    \item 体教融合，开展多样化体育竞赛活动
    \item 减少静态行为，课间"清空教室"，保障课间活动
\end{enumerate}

\textbf{关键目标：}保障学生每天校内体育活动时间不少于1小时，校外体育活动时间不少于1小时，合计每天锻炼时间达到2小时。
\end{keybox}

% ----- 12.2 健康促进学校 -----
\section{健康促进学校}

\subsection{概念}

\begin{defbox}
\textbf{健康促进学校（Health Promoting Schools, HPS）：}通过学校内所有成员（学生、教职员工、家长和社区）的共同努力，制定促进师生健康的规章制度，提供完整、积极的经验和知识结构，包括设置正式和非正式健康课程，创造安全、健康的学校环境，提供适宜的卫生服务，从而促进和保护学生及教职员工健康的学校。

\textbf{核心内涵：}健康促进学校不仅关注健康知识的传授，更强调通过全学校参与的方式，创建有利于健康发展的支持性环境，使健康理念融入学校的日常运作和文化之中。
\end{defbox}

\subsection{全球发展}

\begin{keybox}
\textbf{健康促进学校的全球发展历程：}
\begin{enumerate}[leftmargin=*]
    \item \imp{1995年}：WHO、UNESCO、UNICEF联合发起"全球健康促进学校倡议"（Global School Health Initiative），首次系统提出健康促进学校的概念和框架。
    \item \imp{1990s--2000s}：各国相继开展健康促进学校实践，欧美、澳洲、东南亚等地区广泛推广。
    \item \imp{2021年}：WHO和UNESCO联合发布《全球健康促进学校标准及其指标》（Global Standards for Health Promoting Schools and Their Implementation Guidance），提出8项全球标准和相应的指标体系。
\end{enumerate}
\end{keybox}

\subsection{中国发展}

\begin{keybox}
\textbf{中国健康促进学校发展的三个阶段：}

\begin{enumerate}[leftmargin=*]
    \item \imp{试点阶段（1995年起）}：1995年，中国响应WHO倡导，在北京、上海、天津、浙江等地启动第一批健康促进学校试点项目，初步探索适合中国国情的健康促进学校模式。

    \item \imp{局部推广阶段}：在试点经验基础上，逐步扩大健康促进学校创建范围，多个省份开始区域性推广，形成了以教育部门为主导、卫生部门提供技术支持的协作机制。

    \item \imp{全国推广阶段（2014年至今）}：
    \begin{itemize}
        \item 2014年，国家将健康促进学校创建纳入全民健康素养促进行动。
        \item 2016年，\textbf{《"健康中国2030"规划纲要》}提出"将健康教育纳入国民教育体系，把健康教育作为所有教育阶段素质教育的重要内容"。
        \item 2019年，\textbf{《健康中国行动（2019--2030年）》}之"中小学健康促进行动"明确健康促进学校建设要求。
        \item 2022年，教育部等五部门印发\textbf{《关于全面加强和改进新时代学校卫生与健康教育工作的意见》}，推进全国健康学校建设。
    \end{itemize}
\end{enumerate}
\end{keybox}

\subsection{健康促进学校标准}

\begin{keybox}
\textbf{全球标准（WHO \& UNESCO, 2021）——8项标准：}
\begin{enumerate}[leftmargin=*]
    \item 政府政策与资源支持
    \item 学校政策与资源
    \item 学校治理与领导力
    \item 学校与社区伙伴关系
    \item 学校健康教育课程
    \item 学校社会环境
    \item 学校物质环境
    \item 学校卫生服务
\end{enumerate}
\end{keybox}

\begin{keybox}
\textbf{全学校参与方式（Whole-School Approach）：}
\begin{enumerate}[leftmargin=*]
    \item \imp{学校政策}：将健康纳入学校整体发展规划和政策制度
    \item \imp{课程与教学}：系统开展健康教育，融入学科教学
    \item \imp{学校物质环境}：提供安全、清洁、适宜的校园环境
    \item \imp{学校社会环境}：营造尊重、包容、支持的心理氛围
    \item \imp{学校-社区-家庭联系}：建立家校社协同健康促进网络
    \item \imp{学校卫生服务}：提供综合性学校卫生保健服务
\end{enumerate}
\end{keybox}

\begin{keybox}
\textbf{我国健康学校三大任务目标：}
\begin{enumerate}[leftmargin=*]
    \item \imp{建立健康监测评价机制}：建立健全学校卫生与健康监测体系，对学生健康状况进行动态监测和科学评价，为健康学校建设提供数据支撑。
    \item \imp{增强校园健康服务能力}：完善校医和保健教师配备，提升学校卫生专业服务水平，保障学生获得基本的学校卫生保健服务。
    \item \imp{营造学生健康成长环境}：从课程教学、校园环境、卫生服务、家校社协同等多维度营造促进学生全面健康成长的支持性环境。
\end{enumerate}
\end{keybox}

% ----- 12.3 学校卫生监督 -----
\section{学校卫生监督}

\subsection{内涵}

\begin{defbox}
\textbf{学校卫生监督（school health inspection/supervision）：}卫生行政部门及其卫生监督机构依据国家相关法律法规、卫生标准和规范，对学校（包括托幼机构）的卫生工作进行的监督检查、指导和服务活动。其目的是督促学校改善卫生条件，保障学生健康，预防和控制疾病，促进学生身心健康发展。

\textbf{性质：}学校卫生监督是卫生行政执法行为，具有法定性、规范性、强制性和服务性特征。
\end{defbox}

\subsection{法律依据}

\begin{keybox}
\textbf{学校卫生监督的主要法律依据：}
\begin{enumerate}[leftmargin=*]
    \item \imp{《中华人民共和国宪法》}——国家发展教育事业，保障公民受教育的权利；国家发展医疗卫生事业，保护人民健康
    \item \imp{《中华人民共和国义务教育法》}——保障适龄儿童少年接受义务教育的权利
    \item \imp{《中华人民共和国未成年人保护法》}——保护未成年人的身心健康和合法权益
    \item \imp{《中华人民共和国传染病防治法》}——学校传染病防控的法律依据
    \item \imp{《学校卫生工作条例》（1990年）}——学校卫生监督管理的专门行政法规，是学校卫生监督最直接的法规依据
    \item \imp{《学校卫生监督工作规范》}——具体规范学校卫生监督工作程序和要求
    \item \imp{相关卫生标准与技术规范}：如《中小学校设计规范》（GB 50099）、《中小学校教室采光和照明卫生标准》（GB 7793）、《学校课桌椅功能尺寸及技术要求》（GB/T 3976）、《中小学校传染病预防控制工作管理规范》（GB 28932）等
\end{enumerate}
\end{keybox}

\subsection{主要职责}

\begin{keybox}
\textbf{学校卫生监督的八大职责：}

\begin{enumerate}[leftmargin=*]
    \item \imp{教学及生活环境监督}：对教室（采光照明、微小气候、噪声等）、宿舍、食堂、厕所、浴室等场所的卫生状况进行监督检查，确保符合国家卫生标准。

    \item \imp{传染病防控监督}：监督检查学校传染病预防控制措施落实情况，包括晨午检、因病缺勤追踪、疫情报告、预防接种证查验等制度执行情况。

    \item \imp{饮用水卫生监督}：对学校自备水源、二次供水、直饮水设施、桶装水等的卫生状况和安全管理进行监督检查，确保饮用水符合卫生标准。

    \item \imp{学校医疗机构（保健室）监督}：对校医院、卫生室（保健室）的设置、人员配备、执业许可、药品和医疗器械管理等依法进行监督检查。

    \item \imp{公共场所卫生监督}：对学校内的公共浴室、游泳馆、体育馆、图书馆等公共场所的卫生许可和卫生状况进行监督。

    \item \imp{学校卫生工作协作监督}：与教育行政部门协调配合，共同推进学校卫生工作，落实协作监督机制。

    \item \imp{学校预防性卫生监督}：对新建、改建、扩建校舍的选址、设计、竣工验收进行预防性卫生审查，把好学校卫生第一关。

    \item \imp{其他学校卫生监督}：包括学校突发公共卫生事件应急处置监督、校内食品销售监督等。
\end{enumerate}
\end{keybox}

\subsection{内容方法与评价}

\begin{defbox}
\textbf{学校卫生监督内容——按监督时点分为两大类：}

\textbf{第一大类：预防性卫生监督（preventive supervision）}
\begin{itemize}[leftmargin=*]
    \item \imp{定义}：在学校建设项目的规划、设计、施工和竣工验收阶段进行的卫生审查和监督活动，目的是从源头上保障学校建筑和设施的卫生安全，是学校卫生的第一道防线。
    \item \imp{适用情形}：新建、改建、扩建校舍的选址与设计审查；教学和生活设施设备的配置审查。
\end{itemize}

\textbf{第二大类：经常性/常规性卫生监督（routine supervision）}
\begin{itemize}[leftmargin=*]
    \item \imp{定义}：卫生监督机构按照国家年度监督计划或根据工作需要，对学校的日常卫生工作进行的定期或不定期监督检查。
    \item \imp{主要内容}：教学和生活环境卫生、传染病防控、饮用水卫生、学校医疗机构/保健室、公共场所卫生等。
\end{itemize}

\textbf{监督评价：}根据监督检查结果，对学校卫生工作进行综合评价，出具监督意见书（合格/不合格/限期整改），对违法行为依法实施行政处罚。
\end{defbox}

\subsection{监督程序}

\begin{keybox}
\textbf{学校卫生监督的三大程序：}

\textbf{一、预防性卫生监督审查程序（4个阶段）}
\begin{enumerate}[leftmargin=*]
    \item \imp{项目立项阶段}：审查建设项目的选址和可行性研究报告，提出卫生审查意见
    \item \imp{设计阶段}：审查初步设计和施工图设计中的卫生专篇，确保符合卫生标准
    \item \imp{施工阶段}：开展施工期间的卫生监督检查，确保按审查通过的设计施工
    \item \imp{竣工验收阶段}：参与竣工验收，进行现场卫生学检测，出具卫生学评价意见
\end{enumerate}

\textbf{二、经常性卫生监督检查程序（4个阶段）}
\begin{enumerate}[leftmargin=*]
    \item \imp{准备阶段}：制定监督计划，确定监督对象、内容和人员，准备监督文书和检测设备
    \item \imp{现场检查阶段}：出示执法证件，进行现场检查、采样、检测、询问等，制作现场检查记录
    \item \imp{监督意见阶段}：出具卫生监督意见书，提出整改要求和建议，告知整改期限
    \item \imp{复查与归档阶段}：对整改情况进行复查验收，将监督检查资料整理归档
\end{enumerate}

\textbf{三、行政处罚程序（3个阶段）}
\begin{enumerate}[leftmargin=*]
    \item \imp{立案受理阶段}：发现违法事实后决定立案，指定承办人员，收集固定证据
    \item \imp{调查取证与合议阶段}：全面调查违法事实，告知当事人权利（陈述申辩/听证），经合议形成处罚决定建议
    \item \imp{处罚决定与执行阶段}：下达行政处罚决定书，当事人履行处罚决定，跟踪执行情况，结案归档
\end{enumerate}
\end{keybox}

\newpage
